% Source: https://artofproblemsolving.com/community/c3h2556166_polynomial_question
% Stable identity: aops-topic-2556166-q01
No. Divisibility of polynomials only asks for a polynomial quotient; it says nothing about the coefficients being integers. Examples:
\begin{itemize}
\item $P(x)=x-\dfrac12$ divides $Q(x)=2x^2-x=\left(x-\dfrac12\right)\cdot2x$; here $Q$ has integer coefficients but $P$ does not.
\item $P(x)=x+\sqrt2$ divides $Q(x)=x^2-2=\left(x+\sqrt2\right)\left(x-\sqrt2\right)$.
\item $P(x)=\pi x$ divides $Q(x)=\pi x^2$; neither has integer coefficients.
\end{itemize}
In fact, if $P$ divides $Q$, then so does $cP$ for every nonzero constant $c$ (with quotient $\dfrac{R}{c}$), so a divisor can always be rescaled to have non-integer coefficients. So the answer is $\boxed{\text{no}}$.
